\documentclass[11pt]{article}
\usepackage{amsmath,amssymb,amsthm,hyperref}
\newtheorem{proposition}{Proposition}
\newcommand{\per}{\operatorname{per}}
\title{Two reductions for Marcus's block-permanent inequality}
\author{Research notes for independent verification; no priority claim}
\date{6 November 2026}
\begin{document}
\maketitle
\noindent\textbf{Status.} Neither reduction below proves the full conjecture.
The direct-sum observation identifies the strength a uniform approximate
bound would need. The Gaussian interpolation gives a sufficient inequality
that remains unproved. Equality in Marcus's conjecture is not addressed here.

\paragraph{Direct sums and the published factorial.}
Fix the block size $k$. Write $G(A)_{ij}=\per A_{ij}$ for a positive
semidefinite Hermitian matrix with $m$ blocks. When $\per G(A)>0$, put
$R(A)=\per A/\per G(A)$. For block-aligned orthogonal direct sums,
\[
 G(A\oplus B)=G(A)\oplus G(B),\qquad R(A\oplus B)=R(A)R(B).
\]
Here $k\ge1$, so the permanent of a zero off-diagonal block is zero.

\begin{proposition}
Suppose positive numbers $c_{m,k}$ satisfy
$\per A\ge c_{m,k}\per G(A)$ for every matrix with $m$ blocks of size $k$,
and $\liminf_{m\to\infty}c_{m,k}^{1/m}\ge1$.
Then Marcus's inequality holds for every number of blocks of size $k$.
In particular a positive lower bound depending only on $k$ would suffice.
\end{proposition}
\begin{proof}
For $t$ copies of a fixed matrix with $m$ blocks,
$R(A)^t\ge c_{tm,k}$. Taking $t$th roots and then the lower limit gives
$R(A)\ge1$. If $\per G(A)=0$, the desired inequality follows from
$\per A\ge0$, which follows from its Gram/Fock representation.
\end{proof}

The publisher's abstract of Pate's 1981 paper uses $k$ for the number of
blocks, each of size $n$, and states the factor $1/k!$. Thus in the notation
of these notes the factor is $1/m!$, not $1/k!$. The full text was not
retrieved. This abstract resolves an apparent parameter swap in the 2016
survey and the current Minc-list webpage. The factorial in the number of
blocks does not satisfy the hypothesis above.
Sources: \href{https://doi.org/10.1080/03081088108817400}{Pate (1981), publisher abstract};
\href{https://doi.org/10.1515/spma-2016-0030}{Zhang (2016), p.~313};
\href{https://yair-lavi.com/mincs-list/conjecture-9/index.html}{Minc-list entry}.
The publisher's abstract of
\href{https://doi.org/10.1080/03081088208817427}{Pate (1982)} instead fixes
the number of blocks and lets their size grow, for real symmetric matrices.

\paragraph{The endpoint inequality.}
Use the complex Fock norm, with the inner product conjugate-linear in its
first argument and $\|z^\alpha\|^2=\alpha!$.
Let $p_1,\ldots,p_m$ be arbitrary homogeneous polynomials of one positive
degree $k$ in finitely many variables. Set
\[
 f=\prod_i p_i,\quad f_i=\prod_{j\ne i}p_j,\quad
 G_{ij}=\langle p_i,p_j\rangle,\quad
 B(p)=\sum_{i,j}G_{ij}\langle f_i,f_j\rangle.
\]
The same-index placement in the last expression is intentional. Indeed
\[
 B(p)=\left\|\sum_i p_i(x)f_i(y)\right\|^2,
\]
where $x$ and $y$ are independent sets of standard complex Gaussian
variables. Define $D(p)=m\|f\|^2-B(p)$.
The sufficient inequality to be investigated is
\begin{equation}\label{endpoint}
 D(p)\ge0.
\end{equation}
It is not assumed to follow from previously proved finite cases.

\begin{proposition}
Fix $m,k$. If \eqref{endpoint} holds for every such tuple of arbitrary
homogeneous polynomials in every finite number of variables, then
\[
 \|p_1\cdots p_m\|^2\ge\per G
\]
for all those tuples. In particular it implies the Marcus inequality for
$m$ blocks of size $k$.
\end{proposition}
\begin{proof}
Write $X_i=p_i(Z)$, and let $W$ be an independent centered proper complex
Gaussian vector with $\mathbb E\overline W_iW_j=G_{ij}$. Define
\[
 Y(t)=\sqrt t\,X+\sqrt{1-t}\,W,\qquad
 M(t)=\mathbb E\left|\prod_iY_i(t)\right|^2,\quad 0\le t\le1.
\]
Its covariance is $G$ for every $t$. Put
$Y_{-i}(t)=\prod_{j\ne i}Y_j(t)$ and
$B(t)=\sum_{i,j}G_{ij}\mathbb E\overline{Y_{-i}(t)}Y_{-j}(t)$.
Complex Gaussian integration by parts, or the elementary Gaussian
Ornstein--Uhlenbeck generator applied to $|\prod_i y_i|^2$, gives
\begin{equation}\label{derivative}
 tM'(t)=mM(t)-B(t),\qquad 0<t\le1.
\end{equation}
For clarity, with $s=-\log t$ the drift contributes $-mM$ and the Gaussian
second-derivative term contributes $B$, so $dM/ds=-mM+B$.
All expressions are polynomials in $t$ in this homogeneous circular case,
so the identity also has its endpoint interpretation.

To apply the hypothesis to the mixed-degree vector $Y(t)$, approximate $W$
by
\[
 W_i^{(R)}=R^{-1/2}\sum_{a=1}^R p_i(Z^{(a)}),
\]
using independent copies, also independent of $Z$. For every finite $R$,
$\sqrt t\,p_i(Z)+\sqrt{1-t}\,W_i^{(R)}$ is still a homogeneous degree-$k$
polynomial in finitely many variables. These vectors converge in all the
moments used here to $Y(t)$: expand a joint moment of order at most $2m$
according to copy labels. Terms with a singleton vanish; the surviving
leading terms pair a holomorphic and an antiholomorphic factor and give
the Gaussian Wick moments. Every other nonzero grouping uses fewer copy
labels and tends to zero after normalization. Homogeneity and circular
Gaussian input give zero means, zero unconjugated covariances, and finite
moments of all orders.

Taking limits in \eqref{endpoint} therefore gives $mM(t)-B(t)\ge0$.
Equation \eqref{derivative} implies that $M$ is nondecreasing. Its endpoints
are $M(1)=\|f\|^2$ and $M(0)=\per G$, as required.
For split $p_i$, the two sides are exactly the permanents in Marcus's
conjecture, by the Gram/Fock representation.
\end{proof}

The hypothesis is genuinely dimension-uniform and concerns arbitrary
homogeneous forms. A bound restricted to split forms is insufficient for
this argument, because sums of independent copies need not split.

\paragraph{Combinatorial form.}
For split forms, let $A$ be the Gram matrix of the $mk$ linear factors,
with its designated partition into $m$ blocks. If $b(\sigma)$ counts the
blocks sent wholly onto some block by $\sigma\in S_{mk}$, then exactly
\[
 D(p)=\sum_{\sigma\in S_{mk}}(m-b(\sigma))
                    \prod_{a=1}^{mk}a_{a,\sigma(a)}.
\]
To see this, expand each term of $B$. It chooses one source block and one
target block, matches those blocks internally, and matches their
complements; hence a permutation occurs exactly $b(\sigma)$ times.
The nonnegative integer coefficient does not establish nonnegativity of
the sum, since individual diagonal products can be complex or negative.

\paragraph{Exploration, not a proof of the hypothesis.}
The standalone script \texttt{marcus-gaussian-interpolation.py} minimizes
$D$ by exact one-form Hermitian eigenproblems, with floating-point arithmetic.
It checks its quadratic pencil against direct Fock products and also checks
\eqref{derivative} against finite differences of the full interpolation.
Completed restarts were 100 each for $(m,k,r)=(3,2,4),(4,2,5),(5,2,6),(4,3,4)$,
and 200 for $(3,2,8)$. No negative value exceeding roundoff was found.
Only the derivative at $t=1$ was optimized. Neither this experiment nor the
conditional proposition proves \eqref{endpoint} or the full conjecture.
\end{document}
